\documentclass[11pt]{article}
\usepackage[margin=1in]{geometry}
\usepackage{amsmath,amssymb}
\usepackage[hidelinks]{hyperref}

\title{Weyl group signs do not determine the \(\chi_y\)-genus\\[2pt]
  \large Disproof of Conjecture 00000001131}
\author{}
\date{}

\begin{document}
\maketitle

\noindent
For the ordinary sign distribution, take the two counts of elements with
sign $+1$ and sign $-1$. Two complex semisimple flag varieties have the
same sign distribution, rank, dimension, and Weyl group order, but
different Hirzebruch $\chi_y$-genera.

\section*{The genus and the Weyl length polynomial}

Let $G$ be a connected complex semisimple group, $B$ a Borel subgroup,
and $X=G/B$. Write $W$ for the Weyl group and $\ell$ for its Coxeter
length. The Bruhat decomposition of $X$ has one affine cell of complex
dimension $\ell(w)$ for each $w\in W$~\cite[\S8.3]{springer}. Each such
cell contributes $(uv)^{\ell(w)}$ to the compactly supported
Hodge--Deligne polynomial~\cite[Example~3.5]{mustata}. By additivity and
the smooth projectivity of $X$, this gives
\[
 \sum_{p,q}(-1)^{p+q}h^{p,q}(X)u^pv^q
    =\sum_{w\in W}(uv)^{\ell(w)}.
\]
It follows that $h^{p,q}=0$ for $p\ne q$ and that $h^{p,p}$ counts the
elements of length $p$. The Hirzebruch genus is therefore
\[
 \chi_y(X)=\sum_{p,q}(-1)^q h^{p,q}(X)y^p
          =P_W(-y),\qquad P_W(t)=\sum_{w\in W}t^{\ell(w)}.
\]

\section*{Equal sign counts and unequal genera}

Use groups of types $A_1\times B_3\times G_2\times G_2$ and
$A_2\times A_2\times B_4$, respectively. The Weyl length polynomial is a
product of factors $[d]_t=1+t+\cdots+t^{d-1}$, where $d$ are the basic
invariant degrees~\cite[\S3.15]{humphreys}. The degrees of the groups of
types $A_r$ and $B_r$ are $2,3,\ldots,r+1$ and $2,4,\ldots,2r$, as seen
from symmetric polynomials and symmetric polynomials in the squared
coordinates, respectively. The $G_2$ Weyl group is dihedral of order
twelve, with two alternating elements of each length from one through
five and one element at each of lengths zero and six. Its polynomial is
$[2]_t[6]_t$. Lengths add across direct products, so the degree lists are
\[
  (2,2,2,2,4,6,6,6),\qquad (2,2,2,3,3,4,6,8).
\]
Both lists have eight entries, the rank in each case. Both length
polynomials have degree $22$, the dimension of the largest Bruhat cell,
and each has value $13824$ at $t=1$, the Weyl group order. Multiplication
by a simple reflection pairs elements of opposite sign, so both groups
have $6912$ elements of each sign.

For the coefficient of $t^2$, one can take the linear term from two
distinct factors in $\binom{8}{2}=28$ ways, or take the quadratic term
from one factor with $d\ge3$. There are four such factors in the first
product and five in the second. Thus the coefficient of $t^2$ is $32$
and $33$, respectively. These are also the coefficients of $y^2$ in
the two genera. The identical sign counts cannot determine both values,
even with the rank, dimension, and group order supplied. The polynomial
$\sum_w\operatorname{sgn}(w)y^{\ell(w)}$ does give $\chi_y$, because that
polynomial retains the length of each element in addition to its sign.

\section*{Lean verification}

The Lean 4 project multiplies the eight factors in each degree list by
coefficient convolution. It checks the rank, dimension, group order,
and sign counts, and checks the different $y^2$ coefficients. The
theorem \texttt{C1131.no\_rule\_from\_invariants} rules out any coefficient
rule that uses just the matching data. Run \texttt{lake build} in the
\texttt{lean} directory to verify it.

\begin{thebibliography}{3}
\bibitem{springer}
T.~A. Springer, \emph{Linear Algebraic Groups}, second edition, Birkh\"auser,
1998. \href{https://www.math.utah.edu/~ptrapa/math-library/springer/Springer-Linear-Algebraic%20Groups-1998.pdf}{Book PDF}.

\bibitem{mustata}
M.~Musta\c{t}\u{a}, \emph{Spaces of arcs in birational geometry},
lecture notes. \href{https://websites.umich.edu/mmustata/lectures_arcs.pdf}{Author's copy}.

\bibitem{humphreys}
J.~E. Humphreys, \emph{Reflection Groups and Coxeter Groups}, Cambridge
University Press, 1990.
\href{https://doi.org/10.1017/CBO9780511623646.004}{Chapter 3}.
\end{thebibliography}

\end{document}
